% ------------------------------------------------------------
% Table 1: Disease-level detection
% ------------------------------------------------------------
\begin{table}[t]
\centering
\small
\renewcommand{\arraystretch}{1.10}

\begin{tabular*}{\columnwidth}{@{\extracolsep{\fill}}lcccc@{}}
\toprule
\textbf{Model}
& \textbf{Accuracy}
& \textbf{Macro-F1}
& \textbf{Weighted-F1}
& \textbf{Balanced Acc.} \\
\midrule

Logistic Regression
& \textbf{0.991}
& \textbf{0.989}
& \textbf{0.991}
& \textbf{0.989} \\

XGBoost
& 0.990
& 0.988
& 0.990
& 0.987 \\

Linear SVM
& 0.988
& 0.986
& 0.988
& 0.985 \\

Random Forest
& 0.963
& 0.955
& 0.963
& 0.953 \\

Reference Correlation$^\dagger$
& 0.921
& 0.909
& 0.921
& 0.915 \\

\bottomrule
\end{tabular*}

\vspace{3pt}

\parbox{\columnwidth}{
\footnotesize
$^\dagger$ SingleR-style reference-correlation classifier
(Spearman nearest-centroid mapping), following the template-matching
principle used by reference-mapping methods such as
SingleR \citep{aran2019reference}.
}

\caption{
Disease-level (6-class) detection under patient-independent evaluation.
Simple linear classifiers achieve near-perfect performance, motivating
our decision not to introduce a separate disease-level head in
\ourmethod{}: the full pipeline instead uses Logistic Regression
(bolded above) as a lightweight front-end gate, reserving \ourmethod{}'s
architectural contribution for the substantially harder
lineage$\rightarrow$subtype cascade (Table~\ref{tab:main}).
}
\label{tab:disease}

\end{table}